
\chapter{Descripción de la aplicación}

En este capítulo se presenta la definición de roles e historias de usuario
asociadas a la aplicación propuesta, \texttt{lanbench}, relacionando los
distintos perfiles, sus necesidades y las funcionalidades que el sistema debe
ofrecer. Mediante esta metodología se estructuran los casos de uso principales
y se explicitan las restricciones que permitirán validar el sistema en el
contexto del dataset WebNLG en español. A continuación, se describen la
documentación oficial del proyecto, su arquitectura, los tipos de prueba
automatizada empleados y las tecnologías sobre las que se ha construido.

\section{Metodología}

Las historias de usuario constituyen un instrumento ampliamente utilizado en el
desarrollo de software, en particular en las metodologías ágiles. Su
formulación canónica, popularizada por Cohn~\cite{Cohn2004_UserStories}, sigue
la estructura ``como~$\langle$rol$\rangle$, quiero~$\langle$capacidad$\rangle$,
para~$\langle$beneficio$\rangle$'', y permite alinear cada acción o caso de uso
con su finalidad y con los roles a los que va destinado. Asimismo, las
historias se acompañan de criterios de aceptación que delimitan el alcance
funcional y resultan especialmente útiles para validar el sistema frente a los
requisitos establecidos. El conjunto se contrasta habitualmente con el catálogo
de propiedades INVEST (\textit{Independent, Negotiable, Valuable, Estimable,
	Small, Testable})~\cite{Cohn2004_UserStories}, que actúa como heurística de
calidad para detectar historias demasiado grandes, ambiguas o no comprobables.

Esta metodología se complementa con la documentación de proyecto descrita en la
Sección~\ref{sec:doc-proyecto}, de modo que los criterios de aceptación de las
historias se trasladan a casos de prueba automatizados, manteniendo la
trazabilidad bidireccional entre requisito funcional, implementación y
verificación.

\section{Roles}

El sistema distingue dos planos de rol que se combinan pero nunca se colapsan
en una única etiqueta:

\begin{itemize}
	\item \textbf{Roles de servidor}: definen la capacidad global del usuario sobre la
	      aplicación. Existen únicamente dos: \textit{normal} (por defecto) y
	      \textit{moderator} (elevado, con capacidad para crear datasets, acceder a la
	      API de administración y gestionar criterios).
	\item \textbf{Roles de dataset}: definen la capacidad del usuario sobre un dataset
	      concreto, distinguiendo entre propiedad, anotación, revisión y administración.
\end{itemize}

La Tabla~\ref{tab:roles-permisos} sintetiza la matriz de capacidades de cada
rol funcional sobre las operaciones más relevantes de la aplicación.

\begin{table}[ht]
	\centering
	\footnotesize
	\begin{tabular}{p{5.4cm}cccc}
		\hline
		\textbf{Operación}                                                   & \textbf{Anotador} & \textbf{Revisor} & \textbf{Administrador} & \textbf{Moderador} \\
		\hline
		Registro e inicio de sesión                                          & \checkmark        & \checkmark       & \checkmark             & \checkmark         \\
		Listado de \textit{datasets}                                         & \checkmark        & \checkmark       & \checkmark             & \checkmark         \\
		Creación de \textit{dataset}                                         &                   &                  &                        & \checkmark         \\
		Anotación sobre el \textit{dataset}                                  & \checkmark        &                  & \checkmark$^{*}$       &                    \\
		Revisión sobre el \textit{dataset}                                   &                   & \checkmark       & \checkmark$^{*}$       &                    \\
		Concesión de \textit{permits} a otros usuarios                       &                   &                  & \checkmark             &                    \\
		Gestión de credenciales LLM por \textit{dataset}                     &                   &                  & \checkmark             &                    \\
		Lanzamiento de anotación automática asistida                         &                   &                  & \checkmark             &                    \\
		Consulta de estadísticas globales del \textit{dataset}               &                   &                  & \checkmark             &                    \\
		Consulta de estadísticas personales                                  & \checkmark        & \checkmark       & \checkmark             & \checkmark         \\
		Descarga de XML original                                             & \checkmark        & \checkmark       & \checkmark             & \checkmark         \\
		Descarga de XML extendido (dataset\newline completado)               & \checkmark        & \checkmark       & \checkmark             & \checkmark         \\
		Acceso a la API de administración                                    &                   &                  &                        & \checkmark         \\
		Gestión de criterios de evaluación                                   &                   &                  &                        & \checkmark         \\
		\hline
	\end{tabular}
	\caption{Capacidades de cada rol sobre las operaciones de la aplicación.
		\checkmark indica que el rol dispone de la capacidad; $^{*}$ depende de
		\textit{permits} explícitos sobre el \textit{dataset}.}
	\label{tab:roles-permisos}
\end{table}

A partir de esta separación se derivan los cuatro roles funcionales sobre los
que se articulan las historias de usuario:

\begin{itemize}
	\item \textbf{Anotador}: usuario encargado de transformar tripletas RDF en texto en
	      lenguaje natural y de corregir las traducciones automáticas. Su responsabilidad
	      principal es ampliar la cobertura semántica del corpus.
	\item \textbf{Revisor}: usuario encargado de validar las anotaciones. Constituye una
	      segunda capa orientada a la garantía de calidad, con un trabajo de grano más
	      fino centrado en la evaluación del contexto y del registro lingüístico.
	\item \textbf{Agente IA}: componente automático responsable de la generación de
	      propuestas y de la asistencia a la validación de las anotaciones. La
	      implementación actual contempla proveedores compatibles con OpenAI, así como
	      Groq, Google AI Studio y Anthropic, configurables por dataset mediante
	      credenciales propias.
	\item \textbf{Administrador}: usuario responsable del sistema, capaz de habilitar el
	      acceso a otros usuarios mediante la asignación de roles. Entre sus competencias
	      se incluyen la gestión del dataset y la supervisión de las métricas de calidad.
\end{itemize}

Una vez definidos los roles, se esboza el perfil asociado a cada uno, es decir,
los conocimientos previos necesarios para ejercerlo:

\begin{itemize}
	\item \textbf{Anotador y revisor}: requieren un nivel de español equivalente al C1
	      del MCER y un nivel de inglés equivalente al B2 del MCER; ambos perfiles
	      intervienen como participantes del proceso de \textit{crowdsourcing}. Mientras
	      que el anotador no precisa una formación especializada, la calidad del revisor
	      mejora cuando posee bien un perfil lingüístico, capaz de detectar patrones del
	      lenguaje y de disociarlos por dominio o contexto, bien un perfil de
	      procesamiento del lenguaje natural (PLN), dado que el conocimiento de las
	      dificultades específicas de los modelos en tareas como la generación a partir
	      de tripletas RDF facilita la detección de inconsistencias.
	\item \textbf{Administrador}: orientado tanto a perfiles académicos ---estudiantes de
	      inteligencia artificial, ingeniería ontológica o generación de lenguaje
	      natural--- como a perfiles profesionales ---científicos de datos o lingüistas
	      computacionales---.
\end{itemize}

\section{Decisiones de diseño funcional}
\label{sec:diseno-funcional}

Antes de presentar el catálogo de historias de usuario conviene fijar las
decisiones de diseño funcional que las motivan. Estas decisiones afectan a la
integración del modelo de lenguaje, al control de concurrencia, a la
trazabilidad estadística y a la gobernanza del \textit{dataset}.

\subsection{Integración de modelos de lenguaje en la anotación y revisión}

Las pruebas preliminares realizadas con gestores de modelos ejecutados en
local, en particular Ollama, permitieron constatar limitaciones relevantes en
términos de latencia y de calidad de salida. Estas limitaciones se observaron
al emplear modelos de menor tamaño, como Llama~3 con 7\,000 millones de
parámetros o Qwen con 3\,000 millones de parámetros.

En tales condiciones, incluso tras aplicar una capa intensiva de \textit{prompt
	engineering}, no se obtuvieron resultados suficientemente consistentes para un
flujo de trabajo de anotación y validación de calidad académica. Por ello se
adopta como criterio operativo la utilización preferente de modelos servidos en
la nube, sin renunciar a la posibilidad de ejecución local como opción
secundaria para contextos con mayores restricciones de privacidad o
conectividad.

En este contexto, la alternativa seleccionada es Groq, plataforma de inferencia
en la nube que proporciona acceso mediante API a modelos abiertos de altas
prestaciones y que permite la gestión de claves privadas con caducidad para
sesiones o proyectos concretos. Asimismo, su compatibilidad con interfaces de
tipo OpenAI simplifica la integración dentro de arquitecturas software ya
existentes y reduce el coste de migración tecnológica.

A efectos metodológicos, se recomienda no emplear modelos inferiores a la
familia Llama~3 con 70\,000 millones de parámetros, dado que las pruebas
exploratorias realizadas mostraron una mejora apreciable de la fiabilidad de
las sugerencias automáticas en tareas de corrección y reformulación cuando se
alcanza este umbral.

A partir de estas pruebas, se establece un sistema de clasificación jerárquica
de errores con codificación fija. En particular, el sistema contempla las
siguientes categorías: error ortográfico, error gramatical, error sintáctico,
error de traducción, error de correspondencia con RDF y error de diversidad,
entendido este último como la repetición sustancial de una frase previamente
incluida. Cuando una salida presenta varios defectos simultáneamente, el
sistema asigna únicamente el primer código de error aplicable según el orden
anterior. Junto con dicha clasificación, el modelo propone automáticamente una
alternativa corregida que, de acuerdo con su evaluación, subsana el error
detectado.

No obstante, este mecanismo no se concibe como una autoridad cerrada, sino como
una ayuda supervisada. Por ello, el usuario puede rechazar la clasificación
emitida por el modelo siempre que acompañe dicha acción de un comentario
justificativo en el que explicite el carácter de falso positivo de la
detección. Este requisito introduce trazabilidad sobre las discrepancias entre
humano y modelo y, adicionalmente, permite construir en el futuro un
repositorio de casos conflictivos de gran utilidad para la evaluación
comparativa de distintos LLMs y para el refinamiento posterior del sistema.

Desde el punto de vista funcional, tanto en la creación como en la gestión del
\textit{dataset} puede determinarse el rol operativo del LLM. Se distinguen
tres modalidades. En la primera, denominada \textit{ninguno}, no existe
intervención automática del modelo. En la segunda, denominada
\textit{corrector}, el sistema analiza la frase introducida por el usuario,
aplica el algoritmo de detección jerárquica previamente descrito y propone una
reformulación corregida. En la tercera, denominada \textit{creativo}, el modelo
genera directamente una propuesta de lexicalización procurando evitar los
errores tipificados, de modo que el anotador humano actúa principalmente como
editor o validador posterior. En esta última modalidad no se exige comentario
justificativo cuando el usuario modifica la propuesta, ya que la intervención
humana forma parte natural del proceso creativo y de refinamiento.

\subsection{Concurrencia en anotación y revisión}

La aplicación debe permitir la anotación y la revisión concurrente de un mismo
\textit{dataset} sin que ello produzca duplicidades de trabajo ni una
sobrecarga excesiva en el servidor. Para ello, se introduce el concepto de
\textit{sección}, entendida como un conjunto de entradas reservado
temporalmente en memoria para un usuario cuando este inicia una tarea de
anotación o revisión. Mientras la sesión permanezca activa, dicha sección queda
asociada de forma exclusiva a ese usuario, liberándose automáticamente al
cierre de sesión.

Cada usuario dispone, como máximo, de una única sección asignada por
combinación de \textit{dataset} y rol, distinguiéndose de manera independiente
los roles de anotación y revisión. Si en un momento dado todas las secciones
disponibles de un \textit{dataset} se encuentran ya asignadas, el sistema
informa de esta circunstancia al usuario y evita nuevas reservas. Esta política
permite controlar el acceso concurrente de una forma más eficiente que un
esquema de bloqueo a nivel de entrada individual.

En una primera iteración del sistema, el tamaño de las secciones se fija en
diez elementos. Se trata de una decisión de compromiso entre dos objetivos
contrapuestos. Por una parte, un tamaño elevado reduce la frecuencia de
interacción con el servidor, pero puede disminuir la eficiencia colectiva en
las fases finales de completado del \textit{dataset}, cuando quedan pocas
entradas disponibles. Por otra parte, un tamaño demasiado pequeño incrementa el
coste de gestión en el servidor debido al mayor número de operaciones de
reserva, liberación y reasignación. La elección inicial de un tamaño fijo
responde, por tanto, a un criterio pragmático de simplicidad y rendimiento,
aunque se prevé su parametrización en versiones futuras.

La introducción de secciones se justifica por dos razones fundamentales. En
primer lugar, evita que varios usuarios realicen de forma simultánea el mismo
trabajo sobre las mismas entradas. En segundo lugar, impide que el sistema
tenga que gestionar bloqueos de grano excesivamente fino, esto es, a nivel de
una única entrada, lo que encarecería notablemente el control de concurrencia y
dificultaría la escalabilidad operativa de la plataforma.

\subsection{Estadísticas, trazabilidad y análisis del trabajo}

Con objeto de reforzar la trazabilidad del proceso de construcción del
\textit{dataset}, se incorpora un apartado de estadísticas dentro de la página
de gestión, accesible exclusivamente para usuarios con rol de administrador.
Este módulo permite confirmar la autoría de las tareas de anotación y revisión
realizadas sobre un \textit{dataset} concreto, así como visualizar el tiempo de
trabajo asociado a cada usuario y actividad.

La función principal de este componente no es únicamente descriptiva, sino
también de control metodológico. La identificación de autorías y tiempos de
dedicación facilita auditorías internas, permite analizar la distribución real
de la carga de trabajo y proporciona evidencias empíricas sobre el grado de
participación de cada colaborador en el proceso de curación del corpus.

En fases futuras, este módulo podrá ampliarse con estimaciones analíticas
derivadas de los propios datos de uso generados por la plataforma. En
particular, resulta de interés incorporar métricas desagregadas según tipos de
entrada del \textit{dataset}, complejidad semántica, categoría ontológica o
nivel de intervención correctiva requerida, de modo que pueda caracterizarse
con mayor precisión el coste real de anotación y revisión para distintos
subconjuntos del corpus.

\subsection{Permisos y gobernanza del \textit{dataset}}

La página de gestión incorpora también mecanismos para la concesión de permisos
de anotación, revisión y administración sobre cada \textit{dataset}. Este
diseño responde a una lógica de gobernanza explícita en la que los distintos
niveles de acceso no solo controlan la operatividad de la plataforma, sino que
delimitan responsabilidades diferenciadas dentro del ciclo de vida del recurso.

En la evolución prevista del sistema se contempla, además, la introducción de
un modelo formal de propiedad (\textit{ownership}). Bajo este esquema, el
propietario del \textit{dataset} será el único usuario con capacidad para
eliminarlo, modificar sus propiedades básicas y alterar los permisos de
administración. Esta separación entre propiedad y administración resulta
relevante desde el punto de vista organizativo, ya que permite descentralizar
tareas de gestión cotidiana sin perder un nivel último de control institucional
sobre el recurso.

Los administradores, por su parte, pueden consultar las estadísticas del
\textit{dataset} y conceder permisos de anotación y revisión a otros usuarios.
Adicionalmente, se prevé que en futuras versiones puedan generar URL de
invitación asociadas a un \textit{dataset} y a uno o varios roles concretos.
Tales enlaces se configurarán para un único uso o para múltiples usos, de modo
que se facilite la incorporación controlada de nuevos colaboradores sin
necesidad de asignación manual individual. Este mecanismo resulta especialmente
adecuado para campañas de anotación distribuidas o para entornos docentes y de
investigación con equipos variables.

\section{Historias de usuario}

A continuación se enumeran las historias de usuario agrupadas por bloque
funcional\footnote{La numeración no es estrictamente correlativa: los
	identificadores US-15, US-16, US-18, US-22 y US-23 corresponden a historias
	propuestas durante la elaboración inicial del catálogo que fueron
	posteriormente descartadas o fusionadas con otras. Se conserva la numeración
	original por trazabilidad con la documentación interna del proyecto.}. La
Tabla~\ref{tab:us-catalogo} ofrece un catálogo consolidado del conjunto,
organizado por rol y bloque funcional.

\begin{table}[ht]
	\centering
	\footnotesize
	\begin{tabular}{llp{8.5cm}}
		\hline
		\textbf{Código} & \textbf{Actor}          & \textbf{Título}                                               \\
		\hline
		US-01           & Anotador                & Registro e inicio de sesión                                   \\
		US-02           & Anotador                & Listado de \textit{datasets} disponibles                      \\
		US-03           & Anotador                & Visualización legible de tripletas RDF                        \\
		US-04           & Anotador                & Asignación de secciones de trabajo                            \\
		US-05           & Anotador                & Redacción de sentencias en español a partir de tripletas      \\
		US-06           & Anotador                & Redacción apoyada en la referencia en inglés                  \\
		US-07           & Anotador                & Edición de propuestas generadas automáticamente               \\
		US-08           & Anotador                & Alertas automáticas al finalizar la anotación                 \\
		US-09           & Anotador                & Validación de cobertura RDF                                   \\
		US-10           & Anotador                & Corrección de errores lingüísticos y de referencia            \\
		US-11           & Anotador                & Estadísticas personales del trabajo                           \\
		US-12           & Anotador                & Consulta de errores subsanados en revisión                    \\
		US-13           & Revisor                 & Evaluación humana por criterios de calidad                    \\
		US-14           & Revisor                 & Estadísticas personales del trabajo de revisión               \\
		US-17           & Agente IA               & Generación de alertas por incongruencia                       \\
		US-19           & Administrador           & Subida de \textit{dataset} RDF con opciones declarativas      \\
		US-20           & Administrador           & Descarga de los avances del \textit{dataset}                  \\
		US-21           & Administrador           & Estadísticas por usuario y promedios ponderados               \\
		US-24           & Administrador           & Configuración dinámica de criterios de evaluación (pendiente) \\
		US-25           & Sistema                 & Registro horario de peticiones con redacción de sensibles     \\
		US-26           & Sistema                 & Registro diario de errores 500                                \\
		US-27           & Visitante               & Registro directo como moderador mediante código               \\
		US-28           & Operador                & Generación de lotes de códigos de un solo uso                 \\
		US-29           & Usuario con acceso      & Descarga del XML original del \textit{dataset}                \\
		US-30           & Usuario con acceso      & Descarga del XML extendido (\textit{dataset} completado)      \\
		US-31           & Admin. \textit{dataset} & Gestión cifrada de credenciales LLM por \textit{dataset}      \\
		US-32           & Admin. \textit{dataset} & Nombrado y renombrado del \textit{dataset}                    \\
		US-33           & Admin. \textit{dataset} & Anotación automática asistida por IA                          \\
		US-34           & Admin. \textit{dataset} & Descripción breve opcional del \textit{dataset}               \\
		US-35           & Admin. \textit{dataset} & Selección de modelo de IA del catálogo del proveedor          \\
		\hline
	\end{tabular}
	\caption{Catálogo consolidado de historias de usuario, organizado por
		actor y título. La historia US-24 figura como pendiente: su implementación
		forma parte del trabajo futuro
		(\S\ref{sec:limitaciones-trabajo-futuro}).}
	\label{tab:us-catalogo}
\end{table}

\subsection{Anotador}

\subsubsection{Gestión básica}
\begin{itemize}
	\item \textbf{US-01}: como anotador, quiero registrarme e iniciar sesión para acceder
	      a la plataforma de anotación.
	\item \textbf{US-02}: como anotador, quiero ver el listado de datasets disponibles
	      para seleccionar sobre cuál trabajar.
\end{itemize}

\subsubsection{Visualización de tripletas}
\begin{itemize}
	\item \textbf{US-03}: como anotador, quiero visualizar un conjunto de tripletas RDF
	      de forma natural para comprender la información que debe verbalizarse.
	\item \textbf{US-04}: como anotador, quiero recibir las tripletas agrupadas en
	      secciones de trabajo para que la sesión resulte manejable.
\end{itemize}

\subsubsection{Generación de texto}
\begin{itemize}
	\item \textbf{US-05}: como anotador, quiero redactar sentencias en español a partir
	      de las tripletas RDF para generar datos de entrenamiento.
	\item \textbf{US-06}: como anotador, quiero redactar sentencias en español a partir
	      de la referencia en inglés para apoyarme en ella sin perder el foco en la
	      cobertura RDF.
	\item \textbf{US-07}: como anotador, quiero editar las sentencias generadas
	      automáticamente para corregir errores antes de guardarlas.
	\item \textbf{US-08}: como anotador, quiero recibir alertas automáticas al finalizar
	      una anotación para corregir errores ortográficos, gramaticales y semánticos.
\end{itemize}

\paragraph{Refinamientos de US-08.}
\begin{itemize}
	\item \textbf{US-08-R-01}: las alertas deben cubrir errores ortográficos,
	      gramaticales (sintaxis) y semánticos. En caso de omitirlas, el anotador deberá
	      justificar cada omisión, de modo que el revisor pueda revalidarla.
\end{itemize}

\subsubsection{Validación}
\begin{itemize}
	\item \textbf{US-09}: como anotador, quiero validar si un texto cubre todas las
	      tripletas para garantizar la calidad de la anotación.
	\item \textbf{US-10}: como anotador, quiero corregir errores lingüísticos y de
	      referencia ---falta de variedad lingüística, entidad RDF incorrecta o erratas
	      ortográficas, entre otros--- en los textos generados.
\end{itemize}

\subsubsection{Productividad}
\begin{itemize}
	\item \textbf{US-11}: como anotador, quiero consultar estadísticas de mi trabajo ---por
	      ejemplo, número de anotaciones realizadas--- para conocer mi progreso.
	\item \textbf{US-12}: como anotador, quiero consultar los errores corregidos en
	      revisión para evitar reincidir en ellos.
\end{itemize}

\subsection{Revisor}

\subsubsection{Evaluación humana}
\begin{itemize}
	\item \textbf{US-13}: como revisor, quiero evaluar los textos generados conforme a
	      criterios de calidad ---naturalidad, fluidez, adecuación, completitud, cobertura
	      y diversidad--- para medir la adecuación de la anotación.
\end{itemize}

\paragraph{Refinamientos de US-13.}
\begin{itemize}
	\item \textbf{US-13-R-01}: la evaluación se realiza por frase. Los cinco criterios de
	      frase (naturalidad, fluidez, adecuación, completitud y cobertura) se aplican
	      como un asistente secuencial: hasta que el criterio actual no se decide, el
	      siguiente permanece bloqueado. Se permite retroceder y reescribir una decisión
	      anterior dentro de la misma frase.
	\item \textbf{US-13-R-02}: existe un criterio adicional, la diversidad, decidido una
	      única vez para el conjunto de la entrada por ser comparativo entre frases. Solo
	      se aplica cuando la entrada contiene más de una frase.
	\item \textbf{US-13-R-03}: la decisión es binaria. Una decisión negativa exige
	      introducir un motivo (de longitud máxima 280 caracteres) y, en los criterios de
	      frase, una corrección obligatoria del texto que debe diferir del original.
	\item \textbf{US-13-R-04}: el revisor no puede revisar sus propias anotaciones. La
	      cola de revisión excluye automáticamente cualquier entrada anotada por el
	      usuario que solicita la revisión.
\end{itemize}

\subsubsection{Productividad}
\begin{itemize}
	\item \textbf{US-14}: como revisor, quiero consultar estadísticas de mi propio
	      trabajo ---número de revisiones cerradas y tiempo medio por revisión--- para
	      conocer mi progreso y mi ritmo.
\end{itemize}

\subsection{Agente IA}

\subsubsection{Asistencia inteligente}
\begin{itemize}
	\item \textbf{US-17}: como sistema, quiero generar alertas cuando detecte que una
	      traducción puede no ser válida con respecto al contexto RDF o a la referencia.
\end{itemize}

\subsubsection{Credenciales y modelos por dataset}
\begin{itemize}
	\item \textbf{US-31}: como administrador del dataset, quiero registrar, activar y
	      comprobar mis propias claves de proveedor de IA, de modo que la validación
	      asistida del dataset utilice mi credencial en lugar de la global. Las claves se
	      almacenan cifradas y nunca se devuelven al cliente en claro.
	\item \textbf{US-33}: como administrador del dataset, en un dataset con modo
	      \texttt{generation}, quiero lanzar la anotación automática asistida por IA: el
	      sistema bloquea \textit{N} secciones, las anota una a una en segundo plano y
	      permite reintentar o cancelar tras un fallo del proveedor.
	\item \textbf{US-35}: como administrador del dataset, quiero seleccionar el modelo de
	      IA a partir de la lista pública de modelos que ofrece el proveedor (Groq,
	      Google AI Studio), con un mecanismo de introducción manual como respaldo y
	      mensajes de error claros cuando el catálogo no esté disponible.
\end{itemize}

\subsection{Administrador}

\subsubsection{Gestión de datasets}
\begin{itemize}
	\item \textbf{US-19}: como administrador, quiero subir datasets RDF para utilizarlos
	      en la plataforma. El formulario de creación expone, además del nombre y la
	      descripción, las opciones declarativas \textit{tamaño de sección}, \textit{uso
		      de LLM}, \textit{revisión} y \textit{revisiones adicionales tras corrección}.
	\item \textbf{US-20}: como administrador, quiero descargar los avances del dataset
	      RDF para utilizarlos fuera de la plataforma.
	\item \textbf{US-21}: como administrador del dataset, quiero consultar estadísticas
	      por usuario y un promedio general ponderado de los tiempos de anotación y
	      revisión, junto con la cobertura, los errores subsanados en anotación y los
	      errores subsanados en revisión.
	\item \textbf{US-32}: como administrador del dataset, quiero nombrar un dataset al
	      crearlo (con el nombre del fichero como valor por defecto) y renombrarlo
	      posteriormente desde la página de administración, manteniendo la unicidad del
	      nombre por propietario.
	\item \textbf{US-34}: como administrador del dataset, quiero adjuntar una descripción
	      breve opcional al crear el dataset, de modo que la página de visualización
	      muestre, bajo el nombre, una explicación del contenido.
\end{itemize}

\subsubsection{Configuración}
\begin{itemize}
	\item \textbf{US-24}: como administrador, quiero configurar criterios de evaluación
	      personalizados para el proceso de revisión.
\end{itemize}

\subsection{Operación y auditoría}

\begin{itemize}
	\item \textbf{US-25}: como sistema, quiero registrar las peticiones que incluyan
	      datos en formato JSON o de formulario en ficheros horarios para disponer de
	      trazabilidad operativa. Los campos sensibles, como las contraseñas, se redactan
	      antes de escribirse y la escritura se serializa en un único hilo.
	\item \textbf{US-26}: como sistema, quiero registrar exclusivamente los errores 500
	      en un fichero diario específico para facilitar el diagnóstico técnico, sin
	      incluir respuestas 403 ni 404.
\end{itemize}

\subsection{Registro de moderadores}

\begin{itemize}
	\item \textbf{US-27}: como visitante en posesión de un código de moderador válido,
	      quiero registrarme directamente como moderador para acceder a la superficie de
	      administración sin necesidad de intervención operativa posterior. El código es
	      de uso único y consta de exactamente dieciséis caracteres alfanuméricos.
	\item \textbf{US-28}: como operador, quiero generar lotes de códigos de uso único,
	      mediante un script controlado, para distribuirlos entre quienes deban
	      registrarse como moderadores.
\end{itemize}

\subsection{Descargas desde la visualización del dataset}

\begin{itemize}
	\item \textbf{US-29}: como usuario con acceso a un dataset, quiero descargar el XML
	      original desde la pestaña \textit{Visualización del XML} para inspeccionarlo o
	      reutilizarlo fuera de la plataforma.
	\item \textbf{US-30}: como usuario con acceso a un dataset íntegramente completado,
	      quiero descargar el XML extendido con las anotaciones en español incorporadas,
	      de modo que el corpus anotado pueda reutilizarse externamente.
\end{itemize}

\section{Dependencias entre capacidades}

Las historias de usuario no son independientes: el valor del sistema solo se
manifiesta cuando ciertas capacidades existen previamente. La cadena principal
de dependencias es la siguiente:

\begin{enumerate}
	\item Gestión de usuarios y sesiones.
	\item Control de roles y permisos.
	\item Subida y almacenamiento de datasets.
	\item Acceso al dataset y segmentación en secciones.
	\item Anotación de sentencias.
	\item Validación automática.
	\item Revisión humana.
	\item Estadísticas, exportación y seguimiento.
\end{enumerate}

Sin autenticación no existe trazabilidad por usuario; sin permisos no es
posible repartir el trabajo de manera segura; sin la subida del dataset no hay
base sobre la que anotar; sin segmentación no se puede escalar el trabajo
multiusuario; sin validación no se alcanza el objetivo de calidad del corpus;
sin revisión humana no se cierra el ciclo de aseguramiento de calidad; y sin
registros no se dispone de auditoría suficiente para diagnosticar los errores
del servidor.

\section{Documentación oficial del proyecto}
\label{sec:doc-proyecto}

La documentación funcional, técnica y de planificación se aloja en la carpeta
\texttt{docu\allowbreak menta\allowbreak tion/} del repositorio. Cada documento cumple un propósito
específico y se actualiza siguiendo políticas distintas, lo que garantiza la
trazabilidad entre la intención del producto, su diseño y su evolución a lo
largo del tiempo. Los documentos principales son los siguientes:

\begin{itemize}
	\item \textbf{Historias de usuario}: actúa como referencia funcional única y fuente
	      de verdad para describir qué hace el producto y para quién. Recoge el propósito
	      del proyecto, su alcance, los roles, las historias de usuario con sus criterios
	      de aceptación y las relaciones de dependencia entre capacidades.
	\item \textbf{Diseño técnico}: complementa a las historias de usuario explicando el
	      \textit{cómo}. Contiene el modelo de datos, las claves y restricciones, los
	      procedimientos internos (sesiones activas, asignación de secciones, flujo de
	      anotación y revisión), la convención de nombres de la API y la estructura por
	      capas.
	\item \textbf{Auditorías de arquitectura}: capturan, mediante \textit{snapshots}
	      numerados, el estado no funcional del código (deuda técnica, acoplamientos
	      indebidos, código duplicado, dependencias muertas, desviaciones de estilo). Los
	      archivos previos no se reescriben: cada auditoría posterior arranca con los
	      hallazgos heredados como línea base.
	\item \textbf{Auditorías de cobertura funcional}: registran, historia a historia, el
	      grado de cobertura efectiva del catálogo de historias de usuario, distinguiendo
	      entre las capas implicadas y las pruebas asociadas.
	\item \textbf{Planes de épica y de tarea}: agrupan las historias de usuario en
	      bloques manejables (épicas) y descomponen estas en tareas y subtareas
	      verificables. Su naturaleza es desechable: una vez cerrada la épica, las reglas
	      duraderas que se hayan descubierto durante su ejecución se migran a las
	      historias de usuario o al diseño técnico.
	\item \textbf{Documentos de evaluación de la IA}: recogen los protocolos y los
	      resultados de las pruebas realizadas sobre la calidad de la generación y la
	      corrección asistidas por modelos de lenguaje.
\end{itemize}

A esta documentación de proyecto se añaden los archivos \texttt{CLAUDE.md} y
\texttt{DOCUMENTATION.md}, situados en la raíz del repositorio, que definen,
respectivamente, las normas de colaboración con los agentes que asisten al
desarrollo y el mapa global de toda la documentación.

\section{Arquitectura general del proyecto}

El sistema se ha construido como un monolito modular por capas, ejecutado sobre
Node.js y orquestado mediante contenedores Docker. La separación de
responsabilidades se materializa en una serie de paquetes que se invocan en
cascada desde la frontera HTTP hasta la persistencia, y que se complementan con
utilidades transversales.

\begin{itemize}
	\item \textbf{Rutas} (\texttt{routes/}): exponen los \textit{endpoints} web y de la
	      API. Su responsabilidad se limita a declarar la superficie HTTP y enlazarla con
	      el controlador adecuado, sin contener lógica de negocio.
	\item \textbf{Middlewares} (\texttt{middlewares/}): resuelven cuestiones
	      transversales a las rutas: autenticación, control de acceso por rol, subida de
	      ficheros, registro de peticiones y captura de errores.
	\item \textbf{Controladores} (\texttt{controllers/}): traducen las peticiones HTTP en
	      invocaciones a los servicios, validan el formato de entrada, normalizan la
	      salida y gestionan los códigos de estado. Son adaptadores y no contienen reglas
	      de negocio.
	\item \textbf{Servicios} (\texttt{services/}): concentran los casos de uso de la
	      aplicación. Coordinan repositorios, dominio y utilidades para implementar las
	      reglas funcionales descritas en las historias de usuario.
	\item \textbf{Contratos} (\texttt{contracts/}): definen los DTO de la API y las
	      funciones de proyección que aíslan la representación externa de las entidades
	      internas, evitando filtraciones del modelo persistente hacia el cliente.
	\item \textbf{Repositorios} (\texttt{repositories/}): encapsulan el acceso a la base
	      de datos mediante el cliente de Prisma y exponen operaciones de lectura y
	      escritura coherentes con las invariantes del dominio.
	\item \textbf{Entidades} (\texttt{entities/}) y \textbf{dominio} (\texttt{domain/}):
	      contienen el modelo conceptual y la lógica pura que no depende de la
	      infraestructura.
	\item \textbf{Utilidades} (\texttt{utils/}) y \textbf{constantes}
	      (\texttt{constants/}): proporcionan funciones auxiliares ---lectura y escritura
	      de XML, validadores, clientes de proveedores LLM--- y catálogos de códigos y
	      configuraciones funcionales.
\end{itemize}

La capa de presentación reside en \texttt{public/}, sirviendo las páginas HTML
estáticas que consumen la API. La persistencia se apoya en MariaDB, gestionada
mediante el esquema declarativo de Prisma situado en \texttt{prisma/}. El
conjunto se despliega como tres contenedores Docker ---aplicación, base de datos
y un explorador SQL genérico---, lo que garantiza la reproducibilidad del entorno
de ejecución.

La Figura~\ref{fig:arquitectura-capas} representa gráficamente la cascada
descrita.

\begin{figure}[H]
	\centering
	\resizebox{\linewidth}{!}{%
		\begin{tikzpicture}[
				font=\sffamily\scriptsize,
				main/.style={
						rectangle,
						draw=black,
						fill=gray97,
						minimum width=7.2cm,
						minimum height=0.82cm,
						align=center
					},
				side/.style={
						rectangle,
						draw=black,
						fill=white,
						minimum width=3.7cm,
						minimum height=0.72cm,
						align=center
					},
				stack/.style={
						rectangle,
						draw=black,
						fill=gray75,
						minimum width=2.9cm,
						minimum height=0.72cm,
						align=center
					},
				arrow/.style={-{Latex[length=2.2mm]}, thick},
				support/.style={-{Latex[length=1.8mm]}, dashed}
			]
			\node[main] (browser) {\textbf{Navegador}\\Interfaz HTML consumidora de la API};
			\node[main, below=0.42cm of browser] (public) {\textbf{\texttt{public/}}\\Páginas HTML, CSS y JavaScript servidas por Express};
			\node[main, below=0.42cm of public] (routes) {\textbf{\texttt{routes/}}\\Superficie HTTP web y API};
			\node[main, below=0.42cm of routes] (middlewares) {\textbf{\texttt{middlewares/}}\\Autenticación, roles, subida de ficheros, registro y errores};
			\node[main, below=0.42cm of middlewares] (controllers) {\textbf{\texttt{controllers/}}\\Adaptadores HTTP, validación de entrada y códigos de respuesta};
			\node[main, below=0.42cm of controllers] (services) {\textbf{\texttt{services/}}\\Casos de uso y coordinación de reglas funcionales};
			\node[main, below=0.42cm of services] (repositories) {\textbf{\texttt{repositories/}}\\Acceso persistente mediante cliente Prisma};
			\node[main, below=0.42cm of repositories] (prisma) {\textbf{\texttt{prisma/schema.prisma}}\\Modelo declarativo y cliente ORM};
			\node[main, below=0.42cm of prisma] (mariadb) {\textbf{MariaDB}\\Persistencia relacional};

			\node[side, right=0.75cm of controllers] (contracts) {\textbf{\texttt{contracts/}}\\DTO y proyecciones de API};
			\node[side, right=0.75cm of services] (domain) {\textbf{\texttt{domain/}} / \textbf{\texttt{entities/}}\\Modelo conceptual y lógica pura};
			\node[side, left=0.75cm of services] (utils) {\textbf{\texttt{utils/}} / \textbf{\texttt{constants/}}\\XML, validadores, clientes LLM y catálogos};

			\node[stack, left=0.75cm of routes] (node) {\textbf{Node.js}\\runtime};
			\node[stack, left=0.75cm of middlewares] (express) {\textbf{Express}\\servidor HTTP};
			\node[stack, left=0.75cm of prisma] (orm) {\textbf{Prisma}\\ORM};

			\draw[arrow] (browser) -- (public);
			\draw[arrow] (public) -- (routes);
			\draw[arrow] (routes) -- (middlewares);
			\draw[arrow] (middlewares) -- (controllers);
			\draw[arrow] (controllers) -- (services);
			\draw[arrow] (services) -- (repositories);
			\draw[arrow] (repositories) -- (prisma);
			\draw[arrow] (prisma) -- (mariadb);

			\draw[support] (contracts) -- (controllers);
			\draw[support] (contracts) -- (services);
			\draw[support] (domain) -- (services);
			\draw[support] (utils) -- (services);
			\draw[support] (utils) -- (controllers);
			\draw[support] (node) -- (routes);
			\draw[support] (express) -- (middlewares);
			\draw[support] (orm) -- (prisma);
		\end{tikzpicture}%
	}
	\caption{Arquitectura en capas de \texttt{lanbench}: cascada de rutas,
		\textit{middlewares}, controladores, servicios, contratos, repositorios y
		entidades sobre Node.js, Express y Prisma, con persistencia en MariaDB y
		capa de presentación HTML servida desde \texttt{public/}.}
	\label{fig:arquitectura-capas}
\end{figure}

\section{Modelo de datos}
\label{sec:modelo-datos}

El modelo de datos se declara de forma única en el fichero
\texttt{prisma/schema.\-prisma} y se materializa sobre MariaDB a través del
cliente generado por Prisma. La elección de un \textit{ORM} declarativo
responde a dos requisitos: por un lado, mantener la coherencia entre el modelo
conceptual y el esquema físico mediante un único artefacto versionable; por
otro, ofrecer un cliente tipado que reduzca la superficie de errores en los
repositorios.

Las entidades principales se organizan en cinco grupos funcionales:

\begin{itemize}
	\item \textbf{Usuarios y sesiones}. La entidad \texttt{User} almacena la cuenta, su
	      rol global de servidor y los registros de actividad. La entidad
	      \texttt{ActiveSession} representa la reserva temporal de una sección por parte
	      de un usuario en un modo concreto (anotación o revisión).
	\item \textbf{Datasets}. La entidad \texttt{Dataset} es la raíz funcional del sistema
	      y agrupa los metadatos ---nombre, descripción, idiomas, modo de uso del LLM,
	      tamaño de sección y banderas de revisión---. Cada dataset es responsable de su
	      propio ciclo de vida y de las credenciales LLM asociadas mediante
	      \texttt{Dataset\-Llm\-Credential}.
	\item \textbf{Particionado del trabajo}. Las entidades \texttt{Section} y
	      \texttt{Section\-Assignment} materializan respectivamente la división estática
	      del dataset en bloques de tamaño fijo y la asignación dinámica de cada bloque a
	      un usuario durante una ventana temporal acotada. La separación entre ambas
	      resuelve el control de concurrencia a un grano adecuado, equivalente al
	      concepto de \textit{lease} en sistemas distribuidos.
	\item \textbf{Contenido}. La entidad \texttt{Entry} representa una \textit{entry} del
	      corpus WebNLG y enlaza, mediante relaciones de cardinalidad uno-a-muchos, los
	      conjuntos de tripletas RDF originales y modificados, las lexicalizaciones de
	      referencia y las verbalizaciones en español producidas por la plataforma.
	\item \textbf{Permisos y trazabilidad}. La entidad \texttt{Permit} asocia un usuario
	      y un dataset con uno o varios roles funcionales (\textit{anotador},
	      \textit{revisor}, \textit{administrador}), independientes del rol global. Sobre
	      esta separación se sostienen tanto la gobernanza descrita en el capítulo
	      anterior como la trazabilidad de autorías necesaria para las estadísticas.
\end{itemize}

Esta organización mantiene una correspondencia uno-a-uno entre las invariantes
del dominio (un usuario no puede tener más de una sección activa por modo en un
mismo dataset, un revisor no puede revisar sus propias anotaciones, etc.) y las
restricciones físicas declaradas en el esquema (claves compuestas, índices
únicos y reglas de borrado en cascada controlado). El resultado es un modelo
defensivo: las violaciones funcionales son rechazadas por la propia base de
datos, evitando que dependan exclusivamente de la lógica aplicativa.

\section{Modelo de seguridad y de autenticación}
\label{sec:modelo-seguridad}

El modelo de seguridad de \texttt{lanbench} se articula en cuatro planos
coordinados que cubren autenticación, autorización, protección de secretos y
trazabilidad.

\subsection{Autenticación basada en sesiones}
La plataforma emplea sesiones HTTP gestionadas mediante
\texttt{express-session}, con identificadores opacos almacenados en una
\textit{cookie} con atributos \texttt{HttpOnly} y \texttt{SameSite}. Las
credenciales del usuario se contrastan contra una huella derivada con un
algoritmo de \textit{key stretching} (bcrypt) y no se conservan en claro. Esta
aproximación se considera adecuada para una aplicación monolítica desplegada
bajo un único origen y resulta coherente con las recomendaciones OWASP en
materia de gestión de sesiones~\cite{OWASPASVS2024}.

\subsection{Autorización en capas}
El control de acceso se realiza en dos niveles concatenados. Los
\textit{middlewares} de rol global comprueban el rol de servidor del usuario
antes de exponer la superficie de administración. A continuación, los
\textit{middlewares} de rol por dataset consultan la entidad \texttt{Permit}
para autorizar la operación sobre un dataset concreto. Esta concatenación
responde al principio de mínimo privilegio: ningún usuario obtiene capacidades
sobre un recurso por el mero hecho de poseer un rol global.

\subsection{Gestión cifrada de credenciales LLM}
Las claves de los proveedores de inteligencia artificial se almacenan
exclusivamente en su forma cifrada en el campo \texttt{api\_key\_cipher} de
\texttt{Dataset\-Llm\-Credential}. Una representación abreviada de los últimos
cuatro caracteres (\texttt{key\_last4}) se conserva para confirmación visual
sin reexponer la clave completa. Las claves nunca se devuelven al cliente en
claro, en línea con la historia US-31. La rotación de claves se traduce, en
este modelo, en una sobrescritura de la entidad sin necesidad de migraciones
adicionales.

\subsection{Trazabilidad operativa}
La capa de \textit{middlewares} produce dos registros complementarios. El
primero, descrito en US-25, registra de forma horaria todas las peticiones que
incluyen datos serializados (JSON o formulario), redactando previamente los
campos sensibles ---contraseñas, tokens, claves--- mediante una lista negra
explícita. El segundo, descrito en US-26, persiste únicamente los errores 500
en un fichero diario, excluyendo las respuestas 403 y 404 para evitar el ruido
habitual de \textit{scanners} automatizados. Ambos ficheros se serializan a
través de un único hilo de escritura, garantizando que no se produzca
corrupción por concurrencia.

\subsection{Superficie residual}
Las amenazas residuales contempladas son las clásicas del Top~10
OWASP~\cite{OWASPTop10_2021}: inyección a través de los campos textuales del
editor de anotaciones, secuestro de sesión por exposición accidental del
identificador y desviaciones de configuración entre entornos. Las dos primeras
se mitigan mediante el uso de Prisma como capa de parametrización de consultas
(que impide la concatenación directa de cadenas en SQL) y la activación de la
\textit{cookie} segura en producción. La tercera se controla mediante el
archivo \texttt{.env.example} y la separación estricta entre claves de
desarrollo y de producción.

\section{Tipos de pruebas automatizadas}

La calidad del sistema se sostiene sobre una batería de pruebas automatizadas
estructurada en dos niveles, complementada por una verificación específica del
esquema de base de datos.

\begin{itemize}
	\item \textbf{Pruebas unitarias} (626 casos en la \textit{snapshot} v0.5):
	      ejercitan módulos individuales ---controladores, servicios, repositorios,
	      \textit{middlewares}, entidades, mapeadores DTO y funciones auxiliares--- de
	      forma aislada. Los límites externos, como la base de datos o el transporte
	      HTTP, se sustituyen por dobles de prueba, lo que permite ejecutarlas en
	      milisegundos y mantener un ciclo de realimentación corto durante el
	      desarrollo.
	\item \textbf{Pruebas de integración} (46 casos en la \textit{snapshot} v0.5):
	      recorren flujos completos de principio a fin ---registro y autenticación,
	      importación de un dataset, asistente de revisión, entre otros--- invocando al
	      cliente real de Prisma y tratando la aplicación Express como una caja negra.
	      Su objetivo es detectar incompatibilidades entre componentes que las pruebas
	      unitarias no pueden observar por construcción.
	\item \textbf{Verificación del esquema de base de datos}: un script específico
	      contrasta el esquema activo en MariaDB con el modelo declarado en Prisma. Se
	      ejecuta de forma independiente del marco de pruebas funcional y constituye una
	      salvaguarda tras cada actualización del esquema.
\end{itemize}

En conjunto, la \textit{snapshot} v0.5 del repositorio reúne
\textbf{672 casos de prueba en verde}, según se resume en la
Tabla~\ref{tab:tests-por-tipo}. El inventario completo ---con el fichero, el
nombre humano de cada caso y la historia de usuario que cubre--- se mantiene
en el documento \texttt{documentation/TESTS.md} del repositorio, que actúa
como catálogo canónico de la suite y reemplaza al \textit{placeholder}
T-TRAZ-01 originalmente previsto en
\texttt{doc-process-memory/MEMORY-PROBLEMS-1.md} (§8.2). Cada prueba se
asocia, cuando aplica, a la historia de usuario que cubre, manteniendo así
la trazabilidad funcional.

\begin{table}[ht]
	\centering
	\begin{tabular}{lr}
		\hline
		\textbf{Categoría}                                              & \textbf{Nº de casos} \\
		\hline
		\multicolumn{2}{l}{\textit{Pruebas unitarias}}                                         \\
		\quad Entidades                                                 & 14                   \\
		\quad Autenticación y sesiones                                  & 23                   \\
		\quad Usuarios                                                  & 30                   \\
		\quad \textit{Datasets}                                         & 187                  \\
		\quad Anotaciones                                               & 55                   \\
		\quad Revisiones                                                & 107                  \\
		\quad Administración                                            & 16                   \\
		\quad Estadísticas personales (\textit{Me})                     & 16                   \\
		\quad Reglas lingüísticas en español                            & 43                   \\
		\quad Ollama y clientes LLM                                     & 48                   \\
		\quad XML                                                       & 50                   \\
		\quad Herramientas de operación                                 & 8                    \\
		\quad Contratos, \textit{mappers} y utilidades compartidas      & 150                  \\
		\hline
		\textbf{Subtotal de pruebas unitarias}                          & \textbf{626}         \\
		Pruebas de integración                                          & 46                   \\
		\hline
		\textbf{Total}                                                  & \textbf{672}         \\
		\hline
	\end{tabular}
	\caption[Distribución de los casos de prueba automatizados por
		subsistema según la \textit{snapshot} v0.5 del catálogo canónico.]{Distribución
		de los casos de prueba automatizados por subsistema,
		según la \textit{snapshot} v0.5 del catálogo canónico mantenido en
		\texttt{documentation/TESTS.md}. Los subtotales por subsistema
		consolidan las tablas de la v0.4 con las altas de la v0.5; refactorizaciones
		intermedias han renombrado o unificado un número reducido de casos sin
		alterar las cifras agregadas que aporta el documento (626/46/672).}
	\label{tab:tests-por-tipo}
\end{table}

Existen, además, comandos auxiliares para filtrar por nombre, ejecutar un
único archivo o listar el conjunto de pruebas sin ejecutarlas, lo que
facilita el aislamiento de regresiones.

\section{Tecnologías utilizadas}

A continuación se enumeran las tecnologías clave sobre las que descansa la
aplicación, con una breve indicación del uso concreto que se les ha dado en el
proyecto. Salvo donde se indica lo contrario, se trata de proyectos de software
libre con amplia trayectoria en el ecosistema web.

\begin{itemize}
	\item \textbf{JavaScript / Node.js}: lenguaje y entorno de ejecución del servidor. Se
	      ha empleado JavaScript moderno tanto en el código del servidor como en el de
	      los scripts auxiliares y en la lógica del cliente, lo que ha permitido
	      reutilizar utilidades y mantener una única base de conocimiento en el equipo.
	\item \textbf{Express.js}: marco web minimalista para Node.js. Es la base sobre la
	      que se construyen las rutas, los \textit{middlewares} y la composición de la
	      aplicación, gestionando el ciclo de petición y respuesta sobre el que se apoya
	      el resto de capas.
	\item \textbf{express-session}: complemento de Express dedicado a la gestión de
	      sesiones. Persiste el identificador de sesión en el cliente y mantiene el
	      estado autenticado del usuario, condición previa para aplicar la lógica de
	      roles y permisos.
	\item \textbf{Prisma}: \textit{ORM} declarativo. El esquema
	      \texttt{prisma/schema.prisma} actúa como fuente única de verdad del modelo
	      relacional; el cliente generado se utiliza en todos los repositorios para
	      ejecutar operaciones tipadas contra MariaDB.
	\item \textbf{MariaDB}: sistema gestor de bases de datos relacional compatible con
	      MySQL. Almacena los datasets importados, las entradas, las tripletas, las
	      anotaciones, las revisiones y la trazabilidad de usuarios y roles.
	\item \textbf{Docker y Docker Compose}: herramienta de contenedorización y
	      orquestación local. Permite ejecutar la aplicación, la base de datos y un
	      explorador SQL como servicios coordinados, garantizando que el entorno de
	      desarrollo sea reproducible entre máquinas.
	\item \textbf{Multer}: \textit{middleware} de Express especializado en la subida de
	      archivos. Se utiliza para recibir los datasets XML cargados por los
	      administradores antes de su procesamiento.
	\item \textbf{fast-xml-parser}: analizador de XML eficiente para Node.js. Se emplea
	      para transformar los datasets WebNLG en estructuras manejables, así como para
	      reconstruir el XML extendido con las anotaciones en español.
	\item \textbf{Bootstrap}: marco de componentes y rejilla CSS. Aporta una capa visual
	      homogénea a las páginas servidas desde \texttt{public/} y reduce el esfuerzo
	      invertido en estilos personalizados.
	\item \textbf{jQuery}: biblioteca de manipulación del DOM y de peticiones asíncronas.
	      Se utiliza puntualmente en las páginas para simplificar la comunicación con la
	      API y la actualización dinámica de la interfaz.
	\item \textbf{Prisma Studio y Adminer}: exploradores gráficos del modelo relacional
	      empleados de forma complementaria durante el desarrollo para inspeccionar el
	      contenido de la base de datos sin escribir consultas SQL manuales.
	\item \textbf{Mocha}: ejecutor de pruebas. Coordina las pruebas unitarias e
	      integración, organiza los \textit{suites} y produce los informes en consola.
	\item \textbf{Chai}: biblioteca de aserciones que acompaña a Mocha. Aporta un
	      lenguaje expresivo para describir las expectativas en los casos de prueba.
	\item \textbf{testdouble y proxyquire}: utilidades para crear dobles de prueba y para
	      sustituir dependencias en tiempo de ejecución. Permiten aislar el código bajo
	      prueba del resto del sistema.
	\item \textbf{ESLint}: analizador estático de código. Se ha configurado con reglas
	      adicionales orientadas a seguridad y a buenas prácticas para detectar problemas
	      antes de que lleguen al repositorio.
\end{itemize}
